% GENERATED, edit the French unit or the translation table
\mxunit{automatismes-ch11}{Synthesis problems, reasoning and starter routines}

\begin{fluencysheet}{Tue~27~Apr.~2027}{Synthesis problems across several areas of the curriculum}
\mxfq{Warm-up · Quick}{EQ}{Solve $5x-3=2x+9$.}{$3x=12$, $x=4$.}{}
\mxfq{Core · Problem}{PRO}{A price rises by $\pc{20}$, then a fixed discount of $\num{1000}$ Ar is applied, giving $\num{11000}$ Ar. What was the original price?}{$\num{1,2}p-\num{1000}=\num{11000}$, $p=\dfrac{\num{12000}}{\num{1,2}}=\num{10000}$ Ar.}{}
\mxfq{Core + · Method}{REG}{$f(x)=2x+1$ and $g(x)=-x+7$. Find the point of intersection of the curves, then solve $f(x)<g(x)$.}{$2x+1=-x+7$: $x=2$, point $\pt{2;5}$; $f(x)<g(x)\Leftrightarrow3x<6\Leftrightarrow x<2$.}{}
\mxfq{Review · Quick}{QP}{A recipe for $4$ people needs $250$ g of flour. How much for $6$ people?}{$250\div4\times6=375$ g.}{}
\mxfq{Thread · Quick}{GEO}{A shape of area $4$ cm$^2$ is enlarged with scale factor $3$. What is the area of the enlargement?}{Area $\times3^2$: $36$ cm$^2$.}{Pitfall: multiplying the area by $3$ instead of $3^2$.}
\end{fluencysheet}

\begin{fluencysheet}{Thu~29~Apr.~2027}{Proof by contradiction and by cases}
\mxfq{Warm-up · Quick}{DIV}{What is the parity of $n(n+1)$ for an integer $n$?}{Always even: one of the two consecutive integers is even.}{}
\mxfq{Core · Proof by cases}{DEM}{Show that $n(n+1)$ is even for every integer $n$.}{If $n$ is even, the product contains an even factor. If $n$ is odd, $n+1$ is even. In both cases, the product is even.}{}
\mxfq{Core + · Contradiction}{DEM}{Show by contradiction that the sum of a rational number $r$ and an irrational number $a$ is irrational.}{Suppose $r+a=q\in\Q$; then $a=q-r\in\Q$ (difference of two rationals): contradiction with $a$ being irrational.}{}
\mxfq{Review · Formula}{IMG}{$g(x)=2x+5$. Image of $-1$? Preimage of $11$?}{$g(-1)=3$; $2x+5=11\Leftrightarrow x=3$.}{}
\mxfq{Thread · Formula → graph}{REG}{Plot the points with abscissae $-2\,;\,-1\,;\,0\,;\,1\,;\,2$ on the curve of $f(x)=x^2$; give their ordinates.}{$4\ ;\ 1\ ;\ 0\ ;\ 1\ ;\ 4$: the points are symmetric about the vertical axis.}{}
\end{fluencysheet}

\begin{fluencysheet}{Mon~3~May~2027}{Logical connectives, contrapositive and counter-example in context}
\mxfq{Warm-up · Quick}{DEM}{Write the negation of ``$x>2$'' and then of ``$x>2$ or $x<-1$''.}{$x\leqslant2$; $x\leqslant2$ and $x\geqslant-1$, that is $-1\leqslant x\leqslant2$.}{}
\mxfq{Core · Method}{DEM}{Write the contrapositive of ``if $n^2$ is even, then $n$ is even''.}{``If $n$ is odd, then $n^2$ is odd'' (equivalent to the original statement).}{}
\mxfq{Core + · Counter-example}{DEM}{Disprove: ``every prime number is odd''.}{Counter-example: $2$ is prime and even.}{}
\mxfq{Review · True/false}{PNP}{True or false? Justify. The perimeter of a square is proportional to the length of its side.}{True. $P=4c$: constant of proportionality $4$.}{}
\mxfq{Thread · Reading code}{ALG}{What does $h(3)$ return? Write $h(x)$ as an expression.\mxpseudo{def h(x):\\ \mxtab y = x + 2\\ \mxtab return y * y}}{$h(3)=5\times5=25$; $h(x)=(x+2)^2$.}{}
\end{fluencysheet}
